% Related work. Facts from docs/related_work_survey.md (section letters in the comments); claims to say and
% not say from docs/positioning.md.

\section{Related Work}
\label{sec:related}

% J
\textbf{LLMs in self-adaptive systems} have been surveyed as aids for each step of the MAPE-K loop~\cite{TODO-Li2024-LLMs-in-SAS,TODO-Li2025-GenAI-SAS-roadmap}.
Recent work uses LLMs to reason about unanticipated situations within MAPE-K~\cite{TODO-Maia2026-MAPER}, to refine safety operating rules under a grammar~\cite{TODO-SEAMS2026-grammar-constrained-rules}, and to write adaptation managers that are revised from violated constraints observed in test runs~\cite{TODO-Topfer2026-vibe-coding}.
Separately, quantitative verification at run time has long guided adaptation decisions~\cite{TODO-Calinescu2012-CACM-runtime-QV,TODO-Gerasimou2015-EvoChecker}.
To the best of our knowledge, these two lines have not been combined\todo{verify: the survey found no SEAMS 2026 or 2025 paper that checks LLM-generated adaptation policies with PRISM or PCTL}.

% B, B.3
\textbf{LLM-written generalized plans} are programs or policies that solve every instance of a planning domain.
Silver et al.~\cite{TODO-Silver2024-GenPlan-LLM} showed that LLMs can write such programs for PDDL domains when debugged on training tasks, and later work refined the strategy before the code~\cite{TODO-Stein2026-strategy-refinement}, evolved populations of programs~\cite{TODO-Murray2026-GenePlan}, or checked a property of LLM-written heuristics~\cite{TODO-Pereira2026-property-guided}.
Closest to our loop, Prompt, Prove, Patch~\cite{TODO-Drexler2026-PromptProvePatch} proves that an LLM-written general policy reaches the goal on every training task, with choices among accepted transitions resolved adversarially, and returns counterexamples to the LLM.
These approaches target deterministic domains, where running a policy on an instance verifies it, and their generalization to larger instances is tested, except where a hand-written specification supports a proof~\cite{TODO-Stein2026-provably-complete}.
% A
\textbf{LLM output checked by probabilistic model checking} has so far been used in more limited ways.
Gross et al.\ verify the Markov chain induced by an LLM acting as the policy itself, for one seeded run~\cite{TODO-Gross2025-verifying-LLM-decisions}, and use an LLM to repair the unsafe decisions of a learned policy in a single round~\cite{TODO-Gross2024-counterfactual-LLM}.
Other approaches use the model checker qualitatively~\cite{TODO-Lee2025-VeriPlan}, check a model learned from traces at run time~\cite{TODO-Wang2026-Pro2Guard}, or let the LLM write the MDP and the model checker synthesize the strategy~\cite{TODO-Saccon2025-MDP-generation}.
Prior work combined LLM-written per-state policies with PRISM for single grid-world instances~\cite{TODO-predecessor-LLMTrust2026}.
% C
\textbf{Verifier feedback and resampling} have been compared for deterministic verifiers.
LLM-Modulo frameworks place the LLM as a generator inside a loop of sound critics~\cite{TODO-Kambhampati2024-LLM-Modulo}, but Stechly et al.~\cite{TODO-Stechly2025-self-verification} and Olausson et al.~\cite{TODO-Olausson2024-self-repair} found that sampling the LLM again is about as effective as feeding back errors, at matched budgets.
We compare the two at equal cost with a verifier whose output is quantitative and locates where a controller fails.

% I
\textbf{Generalized policies for probabilistic planning} carry a policy learned on small instances to larger ones.
Decision lists in a relational language~\cite{TODO-Yoon2002-inductive-policy,TODO-Fern2006-decision-lists} and neural policies~\cite{TODO-Toyer2020-ASNets,TODO-Garg2020-SymNet} generalize across sizes, but they are evaluated by simulation, without thresholds or verification.
1--2--3--Go!~\cite{TODO-Azeem2025-123Go} learns decision trees from the optimal strategies of small instances of a parametric MDP and computes their value exactly on each larger instance.
Its trees split on raw state variables against constants for a single reachability objective, so they transfer where good decisions do not depend on the instance's parameters.
Our rule lists share the decision-list form of the earlier work, but an LLM writes them over a vocabulary that names instance values, and they are certified against several requirements at once.
We compare against 1--2--3--Go!'s trees, both over raw variables and over our vocabulary, in \autoref{sec:evaluation}.

% F, H
\textbf{Families of MDPs and permissive strategies} give exact guarantees beyond a single strategy for a single model.
Policy trees~\cite{TODO-Andriushchenko2024-PoliciesGrowOnTrees}, satisfiability-modulo model checking~\cite{TODO-Heck2026-SMPMC} and robust policy gradients~\cite{TODO-Galesloot2025-rfPG} synthesize policies that hold on every member of a finite family, and Schnitzer et al.~\cite{TODO-Schnitzer2025-robust-parametric} bound the risk on unseen members of a distribution.
Their members share one state space, whereas the instances of our families differ in size.
Permissive strategies~\cite{TODO-Draeger2015-permissive} and shields~\cite{TODO-Andriushchenko2026-probabilistic-shields} allow sets of actions with guarantees, and tools such as dtControl~\cite{TODO-Ashok2021-dtControl} compress a strategy into a small decision tree, in each case for a single model.
A certified rule list is a sound permissive strategy on each instance, and we make no claim that it is maximally permissive.
